% ============================================================================
% Fichier  : partie6/P06_C19_tribunal.tex : FORN (LPI342)
% Sources : ACPO ; déontologie de l'expertise. Opération TSANAGA.
% ============================================================================

\chapter{La Présentation devant un Tribunal}
\label{chap:tribunal}
\chaptermark{Devant un tribunal}

\epigraph{%
    « Devant le juge, l'expert ne plaide pas~: il explique. »%
}{Principe du cours}

\niveauChapitre{Intermédiaire à expert}{Chapitre~\ref{chap:rapport}.}

\begin{boxobjectifs}
\begin{itemize}
    \item Préparer une présentation orale claire.
    \item Expliquer une preuve technique à un public non technicien.
    \item Répondre aux questions sans sortir de son rôle.
    \item Tenir face à la contradiction.
    \item Rester dans les limites de l'expertise.
\end{itemize}
\end{boxobjectifs}

\bigskip

% ============================================================================
\section*{Question directrice : comment faire comprendre au tribunal ce que montrent les preuves, sans les trahir ?}
% ============================================================================

Le rapport est écrit~; il faut maintenant l'exposer à l'audience, répondre aux
questions, tenir face à un expert adverse. Ce chapitre prépare cette
restitution orale.

\medskip
\textbf{Plan du chapitre.}
\begin{enumerate}
    \item Préparer la présentation (section~\ref{sec:c19-preparer}).
    \item Expliquer sans trahir (section~\ref{sec:c19-expliquer}).
    \item Répondre et tenir (section~\ref{sec:c19-repondre}).
    \item Rester dans son rôle (section~\ref{sec:c19-role}).
\end{enumerate}

\separateur

% ============================================================================
\section{Préparer la Présentation}
\label{sec:c19-preparer}
% ============================================================================

\begin{boxprocedure}
\textbf{Préparer une présentation au tribunal}
\begin{enumerate}
    \item Résumer les faits établis en quelques phrases simples.
    \item Préparer une explication imagée de chaque notion technique.
    \item Anticiper les questions et les objections.
    \item Préparer les annexes (chaîne de possession, condensats).
    \item Distinguer, à l'oral aussi, le certain, le probable, le NEK.
\end{enumerate}
\end{boxprocedure}

% ============================================================================
\section{Expliquer sans Trahir}
\label{sec:c19-expliquer}
% ============================================================================

Expliquer un condensat, une image bit à bit ou une chaîne de possession à un
magistrat demande des images justes, sans simplification trompeuse. On dit
« une empreinte qui change entièrement si un seul caractère change », non « un
code infaillible ». La pédagogie ne doit jamais déformer le fait.

% ============================================================================
\section{Répondre et Tenir}
\label{sec:c19-repondre}
% ============================================================================

\begin{boxscenario}{Opération TSANAGA. La question de la défense}
\textbf{Constatation.} À l'audience, l'avocat de la défense demande~:
« Comment être sûr que le disque analysé est bien celui de mon client~? »

\medskip
\textbf{Analyse} (confiance élevée). M.~BIKORO répond par la chaîne de
possession et le condensat~: le disque scellé, la fiche de scellé, l'empreinte
identique à l'acquisition et à l'audience. La réponse est factuelle, vérifiable,
calme. Il reconnaît aussi, sans détour, la fragilité de la clé USB et le NEK
du second conteneur.
\end{boxscenario}

% ============================================================================
\section{Rester dans son Rôle}
\label{sec:c19-role}
% ============================================================================

\begin{boxalert}
\textbf{Violation typique.} Conclure à la culpabilité du prévenu.

\smallskip
\textbf{Règle~:} l'expert établit des faits techniques et leur portée~; il ne
qualifie pas l'infraction ni la culpabilité, qui relèvent du juge. Sortir de
ce rôle décrédibilise l'expertise.
\end{boxalert}

\begin{boxmetier}{Expert judiciaire}
\textbf{Calme, clair, dans son rôle.} À la barre, l'expert reste factuel, ne
se laisse pas entraîner dans le débat juridique, reconnaît les limites, et
renvoie toujours à ses pièces. La sérénité et l'honnêteté sont ses meilleures
défenses.
\end{boxmetier}

\separateur

\begin{boxresume}
\textbf{Ce qu'il faut retenir de ce chapitre~:}
\begin{itemize}
    \item \textbf{On prépare} une synthèse simple et des explications imagées.
    \item \textbf{On explique sans trahir} le fait.
    \item \textbf{On répond} par les pièces (chaîne, condensat), calmement.
    \item \textbf{On reconnaît} les limites et les NEK.
    \item \textbf{L'expert ne conclut pas} à la culpabilité.
\end{itemize}
\end{boxresume}

\begin{boxinfo}
\textbf{Points de contrôle de l'élève.} Je sais~: préparer une présentation~;
expliquer une notion technique~; répondre par les pièces~; tenir face à la
contradiction~; rester dans mon rôle.
\end{boxinfo}

\bigskip

\begin{boxexo}[title={\ding{45}\quad Exercice~19.1. Préparer \niveaubase}]
Citez trois éléments à préparer avant une audience.
\end{boxexo}

\begin{boxexo}[title={\ding{45}\quad Exercice~19.2. Expliquer \niveaubase}]
Comment expliquer un condensat à un magistrat~?
\end{boxexo}

\begin{boxexo}[title={\ding{45}\quad Exercice~19.3. Le rôle \niveaubase}]
L'expert peut-il conclure à la culpabilité~?
\end{boxexo}

\begin{boxexo}[title={\ding{45}\quad Exercice~19.4. La question \niveaumoyen}]
Répondez à « comment être sûr que c'est bien le disque de mon client~? ».
\end{boxexo}

\begin{boxexo}[title={\ding{45}\quad Exercice~19.5. Simplifier \niveaumoyen}]
Pourquoi une simplification trompeuse est-elle dangereuse~?
\end{boxexo}

\begin{boxexo}[title={\ding{45}\quad Exercice~19.6. Limites \niveaumoyen}]
Pourquoi reconnaître à l'oral la fragilité de la clé USB~?
\end{boxexo}

\begin{boxexo}[title={\ding{45}\quad Exercice~19.7. Contradiction \niveaumoyen}]
Comment tenir face à un expert adverse~?
\end{boxexo}

\begin{boxexo}[title={\ding{45}\quad Exercice~19.8. Sérénité \niveaumoyen}]
Pourquoi le calme est-il une force à la barre~?
\end{boxexo}

\begin{boxexo}[title={\ding{45}\quad Exercice~19.9. Script \niveauexpert}]
Préparez le déroulé de la présentation orale de TSANAGA.
\end{boxexo}

\begin{boxexo}[title={\ding{45}\quad Exercice~19.10. Objections \niveauexpert}]
Listez cinq objections probables et vos réponses factuelles.
\end{boxexo}

\bigskip

\begin{boxinfo}
\textbf{Pour aller plus loin.}
\begin{itemize}
    \item ACPO~\cite{acpo2012}~; ISO/IEC 27042~\cite{iso27042}.
\end{itemize}
\end{boxinfo}

\begin{boxinfo}
\textbf{Transition.} Reste à progresser~: le Chapitre~\ref{chap:mesurer} clôt
le cours sur la mesure de la qualité de l'enquête.
\end{boxinfo}
